% from ega3-1-fr.tex
\oldpage[III]{164}

$S$-$C$-module filtré, $\operatorname{gr}^\bullet(S)$-$C$-module gradué,
$\operatorname{gr}^{\bullet\bullet}(S)$-$C$-module bigradué : 0, 13.6.6.

Morphisme de suites spectrales : 0, 11.1.2.

Nombre géométrique de composantes connexes d'une fibre : III, 4.3.4.

$C$-objet en groupes, $C$-groupe, $C$-anneau, $C$-module : 0, 8.2.3.

Objet des bords, objet des cobords, objet des cocycles, objet des cycles :
0, 11.2.1.

Objet des images universelles : 0, 13.1.1.

Objet gradué associé à un système projectif satisfaisant à (ML) : 0, 13.4.2.

Objet gradué limité inférieurement (supérieurement) : 0, 11.2.1.

Partie propre (sur $\mathfrak Y$) d'un préschéma formel : III, 3.4.1.

Pleine (sous-catégorie) : 0, 8.1.5.

Pleinement fidèle (foncteur) : 0, 8.1.5.

Polynôme de Hilbert : III, 2.5.3.

Propre (morphisme) de préschémas formels : III, 3.4.1.

Représentable (foncteur) : 0, 8.1.8.

Résolution cohomologique, droite, gauche, homologique, injective, libre,
plate, projective : 0, 11.4.1.

Résolution de Cartan-Eilenberg droite, injective : 0, 11.4.2.

Résolution de Cartan-Eilenberg gauche, projective : 0, 11.6.1.

Rétrocompact (ensemble) : 0, 9.1.1.

Strict (système projectif) : 0, 13.4.2.

Suite spectrale dans une catégorie abélienne : 0, 11.1.1.

Suite spectrale faiblement convergente, régulière, corégulière, birégulière :
0, 11.1.3.

Suite spectrale dégénérée : 0, 11.1.6.

Suite spectrale d'un foncteur relativement à un objet filtré : 0, 13.6.4.

Suites spectrales d'un bicomplexe : 0, 11.3.2 et 11.3.5.

Suites spectrales d'hypercohomologie d'un foncteur par rapport à un complexe :
0, 11.4.3.

Suites spectrales d'hyperhomologie d'un foncteur par rapport à un complexe :
0, 11.6.2.

Suites spectrales d'hyperhomologie d'un foncteur par rapport à un bicomplexe :
0, 11.7.2.

Système de coefficients : 0, 11.8.4.

Unibranche (anneau, point) : III, 4.3.6.

Universellement ouvert (morphisme) : III, 4.3.9.
\end{flushleft}

\clearpage
